\documentclass[letterpaper,11pt]{article}
\usepackage[spanish,mexico]{babel}
\usepackage[utf8]{inputenc}
\usepackage{graphicx}
\usepackage{array}
\usepackage{float}
\usepackage{xcolor}
\usepackage{listings}
\usepackage{tocloft}
\renewcommand{\cftsecleader}{\cftdotfill{\cftdotsep}}
\usepackage[
    left=25mm,
    right=25mm,
    top=35mm,
    bottom=30mm,
    headheight=77pt,
    papersize={21.59cm,27.94cm}
]{geometry}
\usepackage{setspace}
\onehalfspacing
\usepackage{parskip}
\usepackage{microtype}
\usepackage{fancyhdr}
\pagestyle{fancy}
\fancyhf{}
\fancyhead[L]{\includegraphics[height=1.5cm]{portada_img/der.png}}
\fancyhead[R]{\includegraphics[height=1.5cm]{portada_img/izq.png}}
\fancyfoot[C]{\thepage}
\renewcommand{\headrulewidth}{0.4pt}
\renewcommand{\footrulewidth}{0pt}
\usepackage{hyperref}
\hypersetup{
    colorlinks=true,
    linkcolor=blue,
    citecolor=blue,
    filecolor=magenta,
    urlcolor=cyan
}
\lstdefinestyle{output}{
    backgroundcolor=\color{black!5},
    basicstyle=\ttfamily\footnotesize,
    numbers=none,
    frame=single,
    rulecolor=\color{black!30},
    breaklines=true,
    showstringspaces=false,
    xleftmargin=0.5cm,
    xrightmargin=0.5cm
}
\newcommand{\captura}[3]{%
\begin{figure}[H]
    \centering
    \includegraphics[width=\textwidth,height=0.61\textheight,keepaspectratio]{\detokenize{image/#1}}
    \caption{#2}\label{#3}
\end{figure}}

\begin{document}
\hypersetup{pageanchor=false}
\thispagestyle{empty} % Sin encabezado ni pie en la portada

% --- ENCABEZADO DE PORTADA (TABLA DE 3 COLUMNAS) ---
\begin{center}
    % La tabla usa 'm' para centrar verticalmente las imágenes y el texto
    \begin{tabular}{m{0.15\textwidth} m{0.60\textwidth} m{0.15\textwidth}}
        % Logo Izquierdo (UNAM)
        \includegraphics[width=2cm]{portada_img/der.png} & 
        
        % Título Central
        \centering\Huge\textbf{Universidad Nacional Autónoma de México} &
        
        % Logo Derecho (Facultad)
        \raggedleft\includegraphics[width=2cm]{portada_img/izq.png}
    \end{tabular}
\end{center}

\vspace{0.3cm}

% --- INFORMACIÓN DEL REPORTE ---
\begin{center}
    {\LARGE \textbf{Facultad de Ingeniería}}\\[1cm]

    {\LARGE \textbf{Alumno:}}\\[0.3cm]
    {\Large Lara Hernandez Angel Husiel - 320060829} \\[1cm]


    {\LARGE \textbf{Arquitectura Cliente-Servidor}} \\[1cm]

	{\LARGE Grupo: 02 - Semestre: 2027-1} \\[1cm]
    
    {\LARGE \textbf{Examen parcial:}}\\[0.5cm]

    {\LARGE {Comunicación Cliente-Servidor.}}\\[1cm]

    {\LARGE \textbf{Profesor:}}\\[0.2cm]
    {\Large MTRO. CARLOS ALBERTO ROMAN ZAMITIZ}\\[1cm]

    {\Large \textbf{Fecha de realización:}}\\[0.2cm]
    {\Large 15 de octubre del 2026}
\end{center}
\newpage
{\hypersetup{linkcolor=black}\tableofcontents}
\thispagestyle{empty}
\newpage
\setcounter{page}{1}
\hypersetup{pageanchor=true}
\singlespacing

\section{Prueba local}
Lo que se modificó fue que el cliente enviara primero mi nombre, leído desde el teclado con \texttt{fgets()}, y después el servidor respondiera. Además de eso, 
el puerto se recibe desde la línea de comandos y se utiliza \texttt{htons()} para convertirlo al orden de bytes de la red.


\begin{lstlisting}[style=output]
./servidor_examen_parcial 4898
./cliente_examen_parcial 127.0.0.1 4898
\end{lstlisting}

\captura{Captura de pantalla 2026-10-15 a la(s) 12.29.03 p.m..png}{Prueba local del cliente y el servidor en la Mac.}{fig:prueba-local}

Lo que vemos aquí es que el servidor, en la terminal inferior, mostró la IP \texttt{127.0.0.1} y el nombre \texttt{Lara Hernandez Angel Husiel}. Por otra parte, el cliente recibió la respuesta \texttt{Registraste exitosamente tu nombre!}. Después de esto, ambos programas terminaron.

\newpage
\section{Conexión al servidor del profesor}
Ahora bien, para esta prueba se ejecutó el cliente apuntando al servidor \texttt{132.248.59.6}, en el puerto 4898. Se escribió mi nombre desde el teclado y se recibió la respuesta de registro exitoso. También se ejecutó \texttt{date} en la misma terminal para mostrar la fecha y hora.
\begin{lstlisting}[style=output]
./cliente_examen_parcial 132.248.59.6 4898
date
\end{lstlisting}

\subsection{Aclaración sobre las dos conexiones:}
Una disculpa, por error me conecté dos veces al servidor y envié mi nombre en ambas ocasiones. Una conexión la realicé desde la red del laboratorio de iOS y la otra desde PC Puma, por lo que se generaron dos registros desde direcciones IP diferentes. Por este motivo incluyo las dos capturas.

\captura{Captura de pantalla 2026-10-15 a la(s) 12.24.44 p.m..png}{Primera conexión al servidor del profesor y salida de date.}{fig:primera-remota}

\newpage
\subsection{Segunda captura de la conexión:}
En esta captura también aparece la conexión al servidor con la respuesta de registro exitoso. Debajo se conserva la prueba local y la salida de \texttt{date} en la misma terminal.

\captura{Captura de pantalla 2026-10-15 a la(s) 12.32.32 p.m..png}{Segunda captura de la conexión al servidor del profesor y salida de date.}{fig:segunda-remota}


\end{document}
